\documentclass[11pt]{article}

\usepackage[margin=0.9in]{geometry}
\usepackage{amsmath,amssymb}
\usepackage{enumitem}
\usepackage{xcolor}
\usepackage[hidelinks]{hyperref}

% shared solution-box machinery (\ifsolutions toggle); instructor build via
% \def\SolutionsBuild{} - see the Makefile exercise-solutions-light target
% Shared exercise/solution machinery, \input by both latent_dynamics_notes.tex
% (via math_preamble.tex) and kalman_filter_exercise.tex. Kept in one file so
% the toggle and box styling are not duplicated.
%
% Single source, student/instructor toggle:
%   Student PDF (default): solution boxes are dropped entirely.
%   Instructor PDF: built with \def\SolutionsBuild{} prepended (see the
%   Makefile *_solutions targets). Do NOT hand-edit the toggle here.
% Boxes are printer-friendly: light tint fills (<= 12%) that stay legible in
% grayscale. Colors are defined here so the file is self-contained.

\usepackage{environ}   % capture-or-discard env bodies (robust inside lists)
\usepackage[most]{tcolorbox}
\definecolor{solnframe}{rgb}{0.10,0.25,0.55} % calm blue for solution boxes
\newif\ifsolutions
\ifdefined\SolutionsBuild\solutionstrue\fi

% Inline micro-exercise, numbered per section, shown in both builds. Optional
% argument tags it, e.g. \begin{exercise}[optional stretch] for the harder tier.
\newtcolorbox[auto counter, number within=section]{exercise}[1][]{%
    breakable, enhanced, sharp corners,
    colback=black!3, colframe=black!55, boxrule=0.5pt,
    left=6pt, right=6pt, top=3pt, bottom=3pt,
    coltitle=black, colbacktitle=black!10, fonttitle=\bfseries,
    title={Exercise~\thetcbcounter\if\relax\detokenize{#1}\relax\else\ (#1)\fi},
}

% Solution box, instructor build only. In the student build the environment
% body is captured by environ and simply not emitted (discarded, not scanned),
% which is robust even inside enumerate/itemize.
\newtcolorbox{solnbox}{%
    breakable, enhanced, sharp corners,
    colback=solnframe!5, colframe=solnframe!70, boxrule=0.5pt,
    left=6pt, right=6pt, top=3pt, bottom=3pt,
    coltitle=solnframe!85!black, colbacktitle=solnframe!12,
    fonttitle=\bfseries, title={Solution},
}
\ifsolutions
    \NewEnviron{solution}{\begin{solnbox}\BODY\end{solnbox}}
\else
    \NewEnviron{solution}{}% discard body in the student build
\fi


\setlength{\parindent}{0pt}
\setlength{\parskip}{0.5em}
\setlist[enumerate]{topsep=0.2em,itemsep=0.4em,leftmargin=1.6em}

\newcommand{\stepbox}[1]{%
  \vspace{0.25em}
  \fcolorbox{black!35}{black!3}{%
    \begin{minipage}{0.96\linewidth}#1\end{minipage}}}

\begin{document}

\begin{center}
{\Large \textbf{The 1D Gaussian update}}\\[0.2em]
{\small combining a prediction with a measurement (${\sim}\,15$ minutes)}
\end{center}

\vspace{0.5em}

You hold a Gaussian belief about a hidden scalar $x$ (a position, a firing rate,
a state), and you get one noisy measurement $y$:
\[
  x \sim \mathcal{N}(\mu_b, P_b)
  \qquad\text{(prior belief)},
  \qquad\qquad
  y \mid x \sim \mathcal{N}(x, R)
  \qquad\text{(measurement)} .
\]
The posterior $p(x\mid y)$ is again Gaussian. Finding its mean and variance is the
single conditional-Gaussian manipulation at the heart of the Kalman filter. That
is the whole exercise. (The lecture notes write the latent as $z$; here $x$ is the
1D hidden state -- see \S\,\emph{The Kalman filter} in the notes.)

\section*{1. Complete the square}

By Bayes' rule the posterior is the prior times the likelihood,
\[
  p(x\mid y) \;\propto\; p(x)\,p(y\mid x)
  \;\propto\;
  \exp\!\left[-\tfrac{1}{2}\left(
      \frac{(x-\mu_b)^2}{P_b} + \frac{(y-x)^2}{R}
  \right)\right] .
\]
The exponent is quadratic in $x$, so the posterior is Gaussian. Expand and collect
the bracket as $a\,x^2 - 2\,b\,x + \text{const}$.

\begin{enumerate}
  \item Read off the two coefficients:
  \[
    a = \rule{4cm}{0.4pt}
    \qquad\qquad
    b = \rule{4cm}{0.4pt}
  \]

  \begin{solution}
  $a = \dfrac{1}{P_b} + \dfrac{1}{R}$ \quad(the sum of precisions), \qquad
  $b = \dfrac{\mu_b}{P_b} + \dfrac{y}{R}$.
  \end{solution}

  \item Completing the square gives $p(x\mid y)=\mathcal{N}(\mu_{\text{new}},
  P_{\text{new}})$ with $P_{\text{new}} = 1/a$ and $\mu_{\text{new}} = b/a$. Write
  both out:
  \[
    \frac{1}{P_{\text{new}}} = \rule{4cm}{0.4pt}
    \qquad\qquad
    \mu_{\text{new}} = \rule{5cm}{0.4pt}
  \]

  \begin{solution}
  $\dfrac{1}{P_{\text{new}}} = \dfrac{1}{P_b} + \dfrac{1}{R}$: \textbf{precisions
  add}. And $\mu_{\text{new}} = \dfrac{\mu_b/P_b + y/R}{1/P_b + 1/R}$, a
  precision-weighted average of the prior mean and the measurement.
  \end{solution}

  \item Define the \emph{gain} $K = \dfrac{P_b}{P_b+R}$. Show the mean and variance
  take the compact update form
  \[
    \mu_{\text{new}} = \mu_b + K\,(y-\mu_b),
    \qquad\qquad
    P_{\text{new}} = (1-K)\,P_b .
  \]

  \begin{solution}
  Multiply numerator and denominator of $\mu_{\text{new}}$ by $\dfrac{P_bR}{P_b+R}$:
  $\mu_{\text{new}} = \dfrac{\mu_b R + y P_b}{P_b+R}
   = \mu_b\dfrac{R}{P_b+R} + y\dfrac{P_b}{P_b+R}
   = \mu_b + K(y-\mu_b)$, using $\dfrac{R}{P_b+R}=1-K$. And
  $P_{\text{new}} = \dfrac{1}{1/P_b+1/R} = \dfrac{P_bR}{P_b+R}
   = \dfrac{R}{P_b+R}P_b = (1-K)P_b$. The quantity $y-\mu_b$ is the
  \emph{innovation} -- the part of the measurement the prior did not predict.
  \end{solution}
\end{enumerate}

\section*{2. What does the gain do?}

Fill in the limits of $K = \dfrac{P_b}{P_b+R}$:
\[
  R\to 0:\ K\to \rule{1.6cm}{0.4pt}
  \quad
  R\to\infty:\ K\to \rule{1.6cm}{0.4pt}
  \quad
  P_b\to 0:\ K\to \rule{1.6cm}{0.4pt}
  \quad
  P_b\to\infty:\ K\to \rule{1.6cm}{0.4pt}
\]

\begin{solution}
$R\to0\Rightarrow K\to1$ (perfect sensor: trust the measurement);
$R\to\infty\Rightarrow K\to0$ (useless sensor: ignore it);
$P_b\to0\Rightarrow K\to0$ (certain prior: ignore the measurement);
$P_b\to\infty\Rightarrow K\to1$ (no prior: trust the measurement). The gain always
leans toward whichever source is more certain.
\end{solution}

\section*{3. One number}

Take $\mu_b = 8$, $P_b = 4$, $y = 10$, $R = 1$. Compute the gain and the update,
and say whether the posterior mean sits closer to the prior ($8$) or the
measurement ($10$), and why.
\[
  K = \rule{2.5cm}{0.4pt}
  \qquad
  \mu_{\text{new}} = \rule{2.5cm}{0.4pt}
  \qquad
  P_{\text{new}} = \rule{2.5cm}{0.4pt}
\]

\begin{solution}
$K=\dfrac{4}{4+1}=0.8$, \;$\mu_{\text{new}}=8+0.8(10-8)=9.6$, \;
$P_{\text{new}}=(1-0.8)\cdot4=0.8$. The prior ($P_b=4$) is more uncertain than the
sensor ($R=1$), so the update leans toward the measurement, landing at $9.6$.
Observing also shrinks the variance, from $4$ to $0.8$.
\end{solution}

\vspace{0.5em}

\stepbox{%
\textbf{This is one Kalman update.}
\[
  K=\frac{P_b}{P_b+R},
  \qquad
  \mu_{\text{new}}=\mu_b+K(y-\mu_b),
  \qquad
  P_{\text{new}}=(1-K)P_b .
\]
Put a \emph{prediction} step in front -- move the belief with a dynamics model,
which inflates its variance -- then apply this update with each new measurement,
and iterate. That loop is the Kalman filter. For the running-average view, a
moving target, and the vector form, see the extended worksheet
(\texttt{kalman\_filter\_exercise}); for the derivations, the \emph{Kalman filter}
and \emph{smoothing} sections of the notes.}

\end{document}
